\section{What would suffice for a quasi-polynomial recognizer}
\label{sec:global-budget}

\subsection{An explicit parameter target for disc discovery}

The logarithm of the bound in Theorem~\ref{thm:cluster-complexity} is, up to
fixed additive polynomial terms,
\[
 r\log t+k\log(Cd^2)+O\bigl(\log t+\log(k+1)+\log(B+1)\bigr).
\]
This expression is more useful than a single claim that cancellation should
be small. It identifies two different geometric questions: how dispersed the
defects are, and how densely their edges meet tetrahedra.

\begin{corollary}[A sufficient discovery budget]\label{cor:discovery-budget}
Suppose the encoded source sizes, coherent-coordinate bit lengths, and
parameter budgets are polynomial in an original input size $n$. If a promised
vertex meridian satisfies
\[
 r\log(n+1)+k\log(d+1)=O\bigl((\log(n+1))^2\bigr),
\]
then the complete cancellation-cluster search on that source is bounded by
$n^{O(\log n)}$ bit operations.
\end{corollary}
\begin{proof}
Use $t\leq n^{O(1)}$ in the displayed logarithmic bound. Since $d\geq1$,
$k\log(Cd^2)=O(k\log(d+1))$ for the fixed positive constant $C$. The
remaining polynomial factors contribute $O(\log n)$ to the logarithm.
\end{proof}

For bounded $d$, the sufficient region includes
$r=O(\log n)$ and $k=O((\log n)^2)$. For
$d\leq(\log n)^A$ it includes
$k=O((\log n)^2/\log\log(n+2))$ with the same $r$ bound.
For $d=\Theta(n)$ it requires $r+k=O(\log n)$ in this estimate. With
bounded $d$ and fixed $r$, even $k=O(\log n)$ yields polynomial discovery.
These are conditional consequences of a proved local algorithm. No source
transformation controlling these parameters on all unknot inputs is supplied.

\subsection{Accounting for branching, failed attempts, and repairs}

A short successful hierarchy does not bound an algorithm that searches many
unsuccessful ones. The following elementary counting lemma supplies a precise
way to state the missing obligation. It applies to an actual execution tree,
including unsuccessful candidate choices and corrective operations.

\begin{lemma}[Branch-budget accounting]\label{lem:branch-budget}
Let a finite rooted execution tree have height at most $H$. For each internal
vertex $v$, let $b(v)\geq1$ be its actual number of children. Suppose every
root-to-leaf path $P$ satisfies
\[
 \sum_{v\in P,\ v\text{ internal}}\log b(v)\leq R.
\]
The tree has at most $e^R$ leaves and at most $(H+1)e^R$ vertices.
If all work charged to any one vertex is at most $W$, total work is at most
$(H+1)e^R W$.
\end{lemma}
\begin{proof}
At every internal vertex choose a child uniformly. A leaf at path $P$ is
reached with probability $\prod_{v\in P}b(v)^{-1}\geq e^{-R}$. Leaf
probabilities sum to one, proving the leaf bound. Every vertex belongs to
some root-to-leaf path, each of length at most $H$, so counting the union of
these paths with multiplicity gives at most $(H+1)$ times the number of
leaves. The last assertion follows by summing the work charges.
\end{proof}

The height condition controls arbitrarily long unary runs. The branch sum
controls actual alternatives, including abandoned choices. A restart or
repair must either appear in this execution tree or have its separate cost
explicitly charged. A potential that decreases only on the final successful
path does not satisfy this requirement.

For example, a complete faithful implementation would meet the stronger
quasi-polynomial target if it proved all of the following jointly:
\begin{enumerate}
\item Every reached source and all transported objects have polynomial
encoded size, and every surface-search call is complete for its required
topological task under the discovery budget above.
\item Every expensive terminal pattern query has a polynomial-size faithful
flat descriptor whose tiles are through, or is handled by another proved
operation with the same cost bound.
\item Bundle extraction, core certification, every remaining saturated
attachment, and curve/disc transport take at most $n^{O(\log n)}$ work per
execution node, including their certificate replay.
\item The complete execution, including failure, cleanup, and restart nodes,
has $H\leq n^{O(1)}$ and $R=O((\log n)^2)$.
\end{enumerate}
The local results in this report replace concrete terms in this list.
Lemma~\ref{lem:branch-budget} then combines the terms. The list is a research
contract, not a theorem that the present recognizer already satisfies it.

\subsection{The residual geometry after free-boundary condensation}

The unsaturated theorem removes all new relations coming from a proper
vertical attachment. Among attached bordered bases, if $\sigma$ are saturated,
there exists a relative decomposition with exactly $\sigma$ additional
two-handles from those bases.
This suggests measuring complexity by a residual attaching system rather than
by numerical genus or the number of product handles. It also gives a sharper
diagnostic: a failure in an unsaturated region should be a source-chart or
core-certification issue, or a resource limit, for the unmarked census. New
relations from these attachments occur in the saturated residue. Marked
transport remains a separate issue even on an unsaturated base.

Relator count alone is insufficient. A single attaching curve can be long,
can interact with many core handles, and must remain embedded with its source
framing. A useful next reduction would output a small triangulated marked core
and compressed boundary curves, together with a proof of their geometric
meaning. The source-parallelity $30t$ bound is linear in the current source;
it supplies neither a polylogarithmic residual core nor short transported
attaching words automatically.

There is already a relevant algorithmic result: Chambers, Lazarus, de Mesmay,
and Parsa decide contractibility of a compressed boundary curve in a
triangulated three-manifold in time
$2^{O(s^3)}\poly(L)$, where $s$ is the triangulation size and $L$ is the
compressed curve description size \cite{compressed-contractibility}.
Thus, if a future certified reduction gave a residual triangulation with
$s=O(\log n)$ and $L\leq n^{O(1)}$, this existing theorem would yield
$n^{O((\log n)^2)}$ for that contractibility task. For
$s=O((\log n)^{2/3})$ it would yield the stronger $n^{O(\log n)}$ bound.
This is an algebraic substitution into a published theorem, not a new
contractibility algorithm. It also concerns one specified task; finding all
needed curves and recognizing the resulting manifold require additional work.

\subsection{Why the pattern result changes the remaining target}

The incoming disconnected-pattern question required control of a potentially
large global permutation order. Componentwise routing removes that order.
Reblocking removes the need for initially aligned partitions whenever every
component meets both interfaces. Consequently a faithful flat through-tile
descriptor is now sufficient; there is no need to force every original tile
to be connected or to have logarithmic interface width.

This enlarges the target class for a geometric producer in a precise way.
The remaining obstructions are components missing an interface, access to
private tile vertices by later geometry, and unrestricted nested reuse.
These are concrete structural properties that can be logged at the native
pattern-transport boundary. Measuring those failures is more informative than
inferring applicability from a large numerical parallelity multiplicity alone.
